\documentclass[11pt,a4paper]{article}
\usepackage{operationsform}
\begin{document}
\OPStitle{SNS・短文コンテンツ原稿}{TPL-DOC-061}{0.2.0}
\OPSsection{1}{基本情報}
\OPSblanklong{案件・文書の識別子}
\OPSblankshort{作成者と連絡先}
\OPSblankshort{作成日と対象版}
\OPSsection{2}{投稿内容}
\OPSblanklong{対象プラットフォーム}
\OPSblanklong{対象者と投稿目的}
\par\Needspace{4\baselineskip}\noindent\textbf{本文・見出し・リンク案}\par\vspace{1mm}
{\fontsize{8.5}{13}\selectfont
\begin{tabularx}{\linewidth}{@{}XXXX@{}}
\textbf{箇所・ID} & \textbf{記述・出典} & \textbf{確認者・状態} & \textbf{留意点}\\[2mm]
\strut & \strut & \strut & \strut\\[5mm]
\strut & \strut & \strut & \strut\\[5mm]
\strut & \strut & \strut & \strut\\[5mm]
\end{tabularx}
}\vspace{2mm}\par
\OPSblanklong{主張の裏付け}
\OPSsection{3}{配信上の確認}
\OPSblanklong{字数と表示の制限}
\OPSblanklong{引用・画像の権利}
\OPSblanklong{誤解を招く省略}
\OPSblanklong{訂正時の連絡経路}
\OPSsection{4}{確認と履歴}
\par\Needspace{4\baselineskip}\noindent\textbf{参照資料と根拠ID}\par\vspace{1mm}
{\fontsize{8.5}{13}\selectfont
\begin{tabularx}{\linewidth}{@{}XXXX@{}}
\textbf{箇所・ID} & \textbf{記述・出典} & \textbf{確認者・状態} & \textbf{留意点}\\[2mm]
\strut & \strut & \strut & \strut\\[5mm]
\strut & \strut & \strut & \strut\\[5mm]
\strut & \strut & \strut & \strut\\[5mm]
\end{tabularx}
}\vspace{2mm}\par
\OPSblanklong{未確認事項・例外}
\par\Needspace{4\baselineskip}\noindent\textbf{点検・改訂の記録}\par\vspace{1mm}
{\fontsize{8.5}{13}\selectfont
\begin{tabularx}{\linewidth}{@{}XXXX@{}}
\textbf{箇所・ID} & \textbf{記述・出典} & \textbf{確認者・状態} & \textbf{留意点}\\[2mm]
\strut & \strut & \strut & \strut\\[5mm]
\strut & \strut & \strut & \strut\\[5mm]
\strut & \strut & \strut & \strut\\[5mm]
\end{tabularx}
}\vspace{2mm}\par
\end{document}
